\documentclass[11pt]{article}
\usepackage[utf8]{inputenc}
\usepackage[T1]{fontenc}
\usepackage{amsmath}
\usepackage{amssymb}
\usepackage{multicol}
\usepackage{geometry}
\usepackage{tikz}
\usepackage{enumitem}
\usepackage{xcolor}
\usepackage{titlesec}

% Layout
\geometry{a4paper, left=1cm, right=1cm, top=0.5cm, bottom=1.2cm}
\setlength{\columnseprule}{0.4pt}

% Títulos
\titleformat{\section}{\normalfont\Large\bfseries\color{blue}}{}{0em}{}
\titleformat{\subsection}{\normalfont\large\bfseries\color{red}}{}{0em}{}

\title{\textcolor{red}{Segundo Trimestre: Matemática
\\Equação do 1º Grau}}
\author{Professor: Jefferson}
\date{}

\begin{document}

\maketitle
\vspace{-1cm}

\begin{multicols}{2}

\section*{Questões Teóricas}

\subsection*{Questão 1}
O que é uma equação do 1º grau?

\textbf{Comentário:}  
É uma equação em que a incógnita aparece elevada ao expoente 1.

\textbf{Resposta:} Uma equação do tipo $ax + b = 0$, com $a \neq 0$.

\subsection*{Questão 2}
Qual é a incógnita na equação $3x + 5 = 11$?

\textbf{Comentário:}  
A incógnita é a letra que representa o valor desconhecido.

\textbf{Resposta:} $x$

\subsection*{Questão 3}
Uma equação do 1º grau pode ter mais de uma solução?

\textbf{Comentário:}  
Equações do 1º grau sempre possuem uma única solução real.

\textbf{Resposta:} Não, possui apenas uma solução.

\subsection*{Questão 4}
Na equação $7x - 4 = 10$, o número $7$ é chamado de quê?

\textbf{Comentário:}  
O número que multiplica a incógnita é chamado de coeficiente.

\textbf{Resposta:} Coeficiente da incógnita.

\subsection*{Questão 5}
O que significa resolver uma equação?

\textbf{Comentário:}  
Resolver é encontrar o valor da incógnita que torna a igualdade verdadeira.

\textbf{Resposta:} Encontrar o valor que satisfaz a equação.

\subsection*{Questão 6}
Em uma equação, o que acontece quando um termo passa de um lado para o outro?

\textbf{Comentário:}  
Ao mudar de lado, o termo muda de sinal.

\textbf{Resposta:} O sinal é invertido.

\subsection*{Questão 7}
A equação $2x = 10$ é do 1º grau?

\textbf{Comentário:}  
A incógnita aparece com expoente 1.

\textbf{Resposta:} Sim.

\subsection*{Questão 8}
Qual é a forma geral de uma equação do 1º grau?

\textbf{Comentário:}  
Existe uma forma padrão usada para representar esse tipo de equação.

\textbf{Resposta:} $ax + b = 0$, com $a \neq 0$.

\subsection*{Questão 9}
O valor $x = 3$ é solução da equação $x + 2 = 5$?

\textbf{Comentário:}  
Substituímos o valor na equação.

\[
3 + 2 = 5
\]

\textbf{Resposta:} Sim.

\subsection*{Questão 10}
Toda equação do 1º grau representa uma situação de igualdade?

\textbf{Comentário:}  
Equações sempre indicam que dois valores são iguais.

\textbf{Resposta:} Sim.

\section*{Questões Aplicadas}

\subsection*{Questão 11}
A soma de um número com 7 é igual a 15. Qual é esse número?

\textbf{Comentário:}  
Seja $x$ o número.
\[
x + 7 = 15
\Rightarrow x = 8
\]

\textbf{Resposta:} 8

\subsection*{Questão 12}
O dobro de um número é 18. Determine esse número.

\textbf{Comentário:}
\[
2x = 18
\Rightarrow x = 9
\]

\textbf{Resposta:} 9

\subsection*{Questão 13}
Um número menos 5 é igual a 12. Qual é o número?

\textbf{Comentário:}
\[
x - 5 = 12
\Rightarrow x = 17
\]

\textbf{Resposta:} 17

\subsection*{Questão 14}
O triplo de um número, somado com 4, é igual a 25. Qual é esse número?

\textbf{Comentário:}
\[
3x + 4 = 25
\Rightarrow 3x = 21
\Rightarrow x = 7
\]

\textbf{Resposta:} 7

\subsection*{Questão 15}
Carlos tinha um certo valor em dinheiro. Após gastar R\$ 20, ficou com R\$ 50. Quanto ele tinha inicialmente?

\textbf{Comentário:}
\[
x - 20 = 50
\Rightarrow x = 70
\]

\textbf{Resposta:} R\$ 70

\subsection*{Questão 16}
A idade de Ana é o dobro da idade de João. Sabendo que João tem 8 anos, qual a idade de Ana?

\textbf{Comentário:}
\[
x = 2 \cdot 8
\]

\textbf{Resposta:} 16 anos

\subsection*{Questão 17}
Um número somado com seu dobro resulta em 27. Qual é esse número?

\textbf{Comentário:}
\[
x + 2x = 27
\Rightarrow 3x = 27
\Rightarrow x = 9
\]

\textbf{Resposta:} 9

\subsection*{Questão 18}
O perímetro de um quadrado é 40 cm. Qual é a medida do lado?

\textbf{Comentário:}
\[
4x = 40
\Rightarrow x = 10
\]

\textbf{Resposta:} 10 cm

\subsection*{Questão 19}
Um táxi cobra R\$ 5 de bandeirada mais R\$ 3 por quilômetro rodado. Se a corrida custou R\$ 29, quantos quilômetros foram percorridos?

\textbf{Comentário:}
\[
5 + 3x = 29
\Rightarrow 3x = 24
\Rightarrow x = 8
\]

\textbf{Resposta:} 8 km

\subsection*{Questão 20}
Em uma sala há o dobro de meninas em relação aos meninos. Sabendo que há 10 meninos, quantas meninas há?

\textbf{Comentário:}
\[
x = 2 \cdot 10
\]

\textbf{Resposta:} 20 meninas

\end{multicols}

\end{document}

